\section{データベースの手引き}\label{sec:db}
この章は，bib ファイル（文献データベース）の書き方の説明です。エントリ型ごとの具体的な書き方と出力は，第\ref{sec:examples}章にあります。

\subsection{bib ファイルの基本}
bib ファイルは UTF-8 のテキストファイルで，文献1件ごとに「エントリ」を書きます。

\begin{bibtemplate}
@<エントリ型>{<キー>,
  <フィールド名> = {<値>},    ※ フィールドごとに「名前 = {値},」と書く
  <フィールド名> = {<値>},
}
\end{bibtemplate}

\begin{itemize}
\item \textbf{エントリ型}（\bibtype{book}，\bibtype{article} など）は，文献の種類です。どれを選ぶかは，次の節で説明します。
\item \textbf{キー}は，本文から引用するときの名前です（\cmd{parencite}\{aoki2015\} の \texttt{aoki2015}）。半角の英数字で，ほかのエントリと重ならないようにします。
\item \textbf{フィールド}は，著者・書名・年などの情報です。名前と使い方は，\ref{sec:fields}節にあります。値は波括弧 \texttt{\{\ \}} で囲みます。
\item 各エントリは，\textbf{単独で完結した正しい書誌}として書きます。原典と共通する項目（訳書の書名や，再録の著者名など）を，省いたり整形したりして書く必要はありません。省略や整形は，スタイルが行います。
\item スタイルが使わないフィールド（\bibfield{abstract}，\bibfield{isbn}，\bibfield{keywords}，\bibfield{file} など）は，無視されます。Zotero などから書き出した bib をそのまま使えます。
\end{itemize}

\subsection{エントリ型の選び方}\label{sec:choose}
\subsubsection{文献の種類とエントリ型}
文献の種類ごとに，使うエントリ型は次のとおりです。

\begin{center}\small
\begin{tabular}{@{}L{0.44\textwidth}L{0.44\textwidth}@{}}
\toprule
\textbf{文献の種類} & \textbf{エントリ型} \\
\midrule
単著・共著の書籍 & \bibtype{book} \\
編者だけで著者のいない本（論文集・講座など） & \bibtype{collection}（\bibtype{book} でもよい） \\
数巻にわたる書籍の全体／その1巻 & \bibtype{mvbook}／\bibtype{book}（\bibfield{volume}） \\
編書・共著書の中の1つの章・論文 & \bibtype{incollection} \\
雑誌・紀要の論文 & \bibtype{article} \\
新聞・商業雑誌の記事 & \bibtype{article}（\bibfield{entrysubtype}） \\
翻訳書・再録・再刊 & もとの型（\bibtype{book}，\bibtype{article}，\bibtype{incollection}）＋\bibfield{related} \\
電子書籍 & \bibtype{book}（\bibfield{origpublisher}，\bibfield{note}） \\
調査報告書・政府刊行物 & \bibtype{report} \\
報告書の中の論文 & \bibtype{incollection}＋\bibfield{crossref} \\
修士論文・博士論文 & \bibtype{thesis} \\
学会報告（口頭） & \bibtype{inproceedings}（\bibfield{booktitle} なし）／\bibtype{unpublished} \\
予稿集の論文 & \bibtype{inproceedings}（\bibfield{booktitle} あり） \\
ウェブページ・ブログ・オンラインデータベース・オンラインで取得した報告書 & \bibtype{online} \\
コンピュータファイル（MRDF） & \bibtype{dataset} \\
\bottomrule
\end{tabular}
\end{center}

専用の書式がないエントリ型は，近い型の書式で組まれます（\ref{sec:types} 節の末尾）。

\subsubsection{\bibtype{article} と \bibtype{incollection} の違い}\label{sec:artcoll}
どちらも「大きな刊行物の中の1つの文章」ですが，\textbf{その入れ物が違います}。入れ物が何かで選んでください。

\begin{keymarglist}
\item[article]
入れ物が\textbf{雑誌・紀要・年報・新聞}などの定期刊行物のとき。入れ物の名前を \bibfield{journaltitle} に書き，巻・号を \bibfield{volume}・\bibfield{number} に書きます。出版社・出版地・編者は書きません（和文誌の発行元だけは，\bibfield{publisher} に書くと出ます）。

\item[incollection]
入れ物が\textbf{書籍}（編書・共著書・論文集・講座の1巻）のとき。入れ物の書名を \bibfield{booktitle} に書き，編者を \bibfield{editor}，出版社を \bibfield{publisher}，ページを \bibfield{pages} に書きます。
\end{keymarglist}

出力の形もはっきり違います。\bibtype{article} は「『誌名』巻(号): ページ」，\bibtype{incollection} は「編者編『書名』出版社，ページ」です。

\bibexample{mori2018}
\bibexample{ono2014}

\begin{caution}
迷いやすいものをまとめます。

\textbf{年報・紀要・学会誌}は，雑誌と同じ \bibtype{article} です（\bibfield{journaltitle} に誌名）。\textbf{講座・叢書の1巻に載った章}は，その巻が書籍なので \bibtype{incollection} です（講座名は，親の巻の \bibfield{series} に書きます）。\textbf{共著書の中の1章}も，\bibtype{incollection} を使います。biblatex には \bibtype{inbook} もありますが，このスタイルでは \bibtype{incollection} と同じ書式で組まれるので，\bibtype{incollection} だけ覚えれば足ります。
\end{caution}

\subsubsection{\bibtype{book}，\bibtype{collection}，\bibtype{mvbook}}
著者のいる本は \bibtype{book}，編者だけで著者のいない本（論文集など）は \bibtype{collection} です。書式は同じで，\bibfield{author} がなく \bibfield{editor} だけのときは，著者の位置に編者が出て「編」が付きます。編書の章を \bibfield{crossref} で指すとき（次の節）は，親を \bibtype{collection} にします。

数巻にわたる本は，全体を \bibtype{mvbook}（\bibfield{volumes} に巻数），1巻だけを \bibtype{book}（\bibfield{volume} に巻の番号）で書きます。

\subsubsection{\bibfield{crossref}：所収書の情報を1か所にまとめる}\label{sec:crossref}
同じ編書の章を何本も引用するときは，所収書の情報（書名・編者・出版社・出版年）を章ごとに繰り返して書くことになり，書き間違いのもとです。\bibfield{crossref} を使うと，所収書を \bibtype{collection} として1回だけ書き，章のエントリには，章に固有の情報（著者・題名・ページ）と \texttt{crossref = \{<親のキー>\}} だけを書けます。

\begin{bibtemplate}
@incollection{<キー>,
  author   = {<姓>,<名>},
  sortname = {<Family>, <Given>},   ※ 和文のみ
  title    = {<論文名>},
  pages    = {<開始ページ>--<終了ページ>},
  crossref = {<親のキー>},          ※ 親は @collection（同じ bib ファイルに書く）
  langid   = {japanese},
}
\end{bibtemplate}

まず，親になる所収書を \bibtype{collection} として書きます（キーは \texttt{sakaied2014}）：

\bibexample{sakaied2014}

章のエントリは，\bibfield{crossref} で親を指します。親の書名・編者・出版社・出版年が引き継がれて，所収書の情報つきで出ます（章だけを引用した場合の出力です。親は出ません）：

\bibexample{onocr2014}

\begin{itemize}
\item 引き継がれるのは，親の \bibfield{title}（子では所収書名 \bibfield{booktitle}），\bibfield{subtitle}（\bibfield{booksubtitle}），\bibfield{editor}，\bibfield{publisher}，\bibfield{location}，\bibfield{date} などです。子に同じフィールドがあれば，子の値が優先されます。
\item 親のエントリは，同じ bib ファイル（または \cmd{addbibresource} で読み込んだファイル）になければなりません。キーを書き間違えると，biber が警告を出します。
\item \textbf{章を引用しても，親は文献リストに載りません。}所収書そのものも載せたいときは，本文で引用するか，\cmd{nocite}\{<親のキー>\} を書きます。
\item 和文の著者（子の \bibfield{author}）には，子にも \bibfield{sortname} が要ります。
\item 親と子の組み合わせは，\bibtype{collection}・\bibtype{book}・\bibtype{proceedings} を親にした \bibtype{incollection}・\bibtype{inproceedings}，\bibtype{report} を親にした \bibtype{incollection}（報告書の中の論文，\ref{sec:report} 節），\bibtype{mvbook} を親にした \bibtype{book}（多巻本の1巻）です。
\end{itemize}

\subsubsection{\bibfield{related} と \bibfield{crossref} の違い}\label{sec:related}
どちらも別のエントリを指しますが，目的が違います。取り違えやすいので，表にまとめます。

\begin{center}\small
\begin{tabular}{@{}L{0.14\textwidth}L{0.36\textwidth}L{0.36\textwidth}@{}}
\toprule
 & \bibfield{crossref} & \bibfield{related}（＋\bibfield{relatedtype}） \\
\midrule
指すもの & 所収書・報告書・多巻本の全体など，文献を\textbf{入れる側} & 原典・初出など，同じ文章の\textbf{別の版} \\
目的 & 書誌情報を引き継ぎ，入力を省く & 出力を1つの項目にまとめる \\
使う場面 & 編書の章，報告書の中の論文，多巻本の1巻 & 訳書（\opt{translationof}），再録・再刊（\opt{reprintof}） \\
指された側 & 引用しなければ文献リストに載らない & 原典が代表として載り，訳書・再録が添えられる \\
\bottomrule
\end{tabular}
\end{center}

「訳書の原典」を \bibfield{crossref} で指したり，「所収書」を \bibfield{related} で指したりしないでください。訳書・再録の書き方は，\ref{sec:trans} 節にあります。

\subsection{エントリ型の一覧}\label{sec:types}
このスタイルが専用の書式を持つエントリ型の，必須・任意のフィールドです。どのエントリ型にも，\bibfield{langid}（\ref{sec:cautions} 節）を書きます。和文の著者・編者には \bibfield{sortname} も書きます。

\begin{typelist}
\typeitem{book}

単著・共著の書籍です。\bibtype{booklet}，\bibtype{manual}，\bibtype{reference}，\bibtype{proceedings}，\bibtype{bookinbook} も同じ書式で組まれます。

\reqitem{author, title, date, publisher, langid}
\optitem{editor, translator, subtitle, edition, series, number, volume, part, maintitle, location, origdate, origpublisher, origlocation, urldate, url, doi, note, related, relatedtype}

\typeitem{collection}

編者だけで著者のいない本です。\bibfield{author} の代わりに \bibfield{editor} を書きます。書式は \bibtype{book} と同じです。

\reqitem{editor, title, date, publisher, langid}
\optitem{subtitle, edition, series, number, location, translator, urldate, url, note}

\typeitem{mvbook}

数巻にわたる書籍の全体です。\bibtype{mvcollection}，\bibtype{mvreference}，\bibtype{mvproceedings} も同じ書式で組まれます。

\reqitem{author, title, date, volumes, publisher, langid}
\optitem{editor, subtitle, location}

\typeitem{incollection}

編書・共著書の中の1つの章・論文です。\bibtype{inbook}，\bibtype{inreference}，\bibtype{suppbook}，\bibtype{suppcollection} も同じ書式で組まれます。

\reqitem{author, title, booktitle, date, publisher, pages, langid}
\optitem{editor, booksubtitle, edition, location, translator, series, number, origdate, urldate, url, note, related, relatedtype, crossref}

\typeitem{article}

雑誌・紀要・年報・新聞の記事です。\bibtype{misc}，\bibtype{software}，\bibtype{patent} も同じ書式で組まれます。

\reqitem{author, title, journaltitle, date, langid}
\optitem{subtitle, volume, number, pages, publisher, howpublished, doi, url, urldate, eprinttype, entrysubtype, issue, note, related, relatedtype}

\typeitem{report}

調査報告書・政府刊行物です。\bibfield{author}，\bibfield{editor}，\bibfield{institution} のうち，先にあるものが先頭の名前になります。

\reqitem{title, date, langid}
\optitem{author, editor, institution, type, number, publisher, location, subtitle, url, urldate}

\typeitem{thesis}

修士論文・博士論文です。

\reqitem{author, title, date, institution, type, langid}
\optitem{subtitle, url, urldate}

\typeitem{inproceedings}

学会報告です。\bibfield{booktitle} があれば，刊行された予稿集の論文として，\bibtype{incollection} と同じ形で出ます。\bibtype{unpublished}（未刊行の報告）も同じ書式で組まれます。

\reqitem{author, title, date, eventtitle, langid}
\optitem{type, eventdate, venue, booktitle, editor, publisher, location, pages}

\typeitem{online}

ウェブページ・ブログ・オンラインデータベース・オンラインで取得した報告書です。

\reqitem{title, date, url, urldate, langid}
\optitem{author, titleaddon, organization, location, subtitle, note}

\typeitem{dataset}

コンピュータファイル（MRDF）です。

\reqitem{author, title, date, type, publisher, langid}
\optitem{location, howpublished}
\end{typelist}

\subsection{フィールドの一覧}\label{sec:fields}
このスタイルが使うフィールドです。（　）の中は，biblatex でのデータの種類で，書き方に影響します（literal＝文字列，name＝名前，date＝日付，range＝ページの範囲，uri＝URL，key＝決まった語，entry set＝別のエントリのキー）。リストの種類のフィールドは，複数の値を \texttt{and} でつなげて書けます。

\subsubsection{名前}
\begin{fieldlist}
\listitem{author}{name}

著者です。複数なら \texttt{and} でつなぎます。和文は \texttt{\{姓,名\}}，欧文は \texttt{\{Family, Given\}}，団体は \texttt{\{\{団体名\}\}} と書きます（\ref{sec:cautions} 節）。

\listitem{editor}{name}

編者です。「編」（欧文は \texttt{ed.}，\texttt{eds.}）が付きます。\bibfield{author} がないときは，著者の位置に出ます。

\listitem{translator}{name}

訳者です。「訳」（欧文は \texttt{trans.}）が付きます。原典に添えた訳書では著者の位置に，単独の訳書では書名の後に出ます。

\fielditem{translatortype}{key}

訳者の役割を「監訳」などに変えるときに書きます。\bibfield{editortype}（編者）も同様です。

\listitem{sortname}{name}

並べ替え用の名前です。和文の著者・編者には，必ず書きます（\ref{sec:cautions} 節）。
\end{fieldlist}

\subsubsection{タイトル・版・巻}
\begin{fieldlist}
\fielditem{title}{literal}

書名・論文名・記事名です。和文では，文脈によって『』または「」で囲まれます。

\fielditem{subtitle}{literal}

副題です。和文は「――」，欧文は「:」で書名につながります。

\fielditem{titleaddon}{literal}

書名に添える語です。\bibtype{online} では，ウェブサイト名・ブログ名を書きます。

\fielditem{booktitle}{literal}

所収書の書名です（\bibtype{incollection}，\bibtype{inproceedings}）。\bibfield{booksubtitle} は副題です。

\fielditem{journaltitle}{literal}

誌名・紙名です（\bibtype{article}）。

\fielditem{edition}{integer または literal}

版です。数字（\texttt{2}）は和文で「第2版」，欧文で \texttt{2nd ed.} になります。「新版」「改訂版」などの語はそのまま出ます。和文では書名に含まれます。

\fielditem{series}{literal}

叢書名・講座名です。和文だけ，書名の前に出ます（『講座社会学 12 環境』）。\bibfield{number} が叢書の中の番号です。

\fielditem{volume}{integer または literal}

多巻本の巻の番号，または雑誌の巻です。\bibfield{part} は巻の中の分冊，\bibfield{volumes} は多巻本の総巻数，\bibfield{maintitle} は多巻本の総書名です。

\fielditem{number}{literal}

雑誌の号，または叢書の中の番号，報告書の課題番号です。
\end{fieldlist}

\subsubsection{日付}
\begin{fieldlist}
\fielditem{date}{date}

刊行年です。\texttt{2015}，または \texttt{1975-03-17}（新聞・雑誌・ブログ）と書きます。年が不明なときは書きません（\texttt{n.d.} と出ます）。刊行中の叢書・多巻本のように終了年がないときは，\texttt{1956/} と書きます（\texttt{1956-} ではありません。biblatex の書き方です）。年は開始年（1956）だけが出ます。

\fielditem{origdate}{date}

原典の刊行年です。訳書と再刊で「原典の年」を出すときと，電子書籍で紙媒体の年を出すときに使います。

\fielditem{urldate}{date}

URL を見た日（取得日）です。これがあるときだけ，（取得日，URL）が出ます。

\fielditem{eventdate}{date}

学会の開催日です（欧文の学会報告で，月日だけが出ます）。

\fielditem{pubstate}{key}

近刊のときは \texttt{forthcoming} と書きます（「近刊」，欧文は \texttt{forthcoming}）。
\end{fieldlist}

\subsubsection{出版・機関・その他}
\begin{fieldlist}
\listitem{publisher}{literal}

出版社です。著者・編者・機関名と同じときと，「国立印刷局」のときは，自動で省かれます。

\listitem{location}{literal}

出版地です。欧文だけ出ます。

\listitem{institution}{literal}

大学・研究機関・研究科などです（\bibtype{thesis}，\bibtype{report}）。

\listitem{organization}{literal}

ウェブサイトの運営機関です（\bibtype{online}）。欧文だけ出ます。

\listitem{origpublisher}{literal}

電子書籍の，紙媒体の出版社です。\bibfield{origlocation} は紙媒体の出版地（欧文だけ）です。

\fielditem{type}{key または literal}

学位論文・報告書・コンピュータファイルの種類です。\texttt{phdthesis}，\texttt{mathesis}，\texttt{techreport} のようなキー，または \texttt{2019年度修士論文} のような文字で書きます。

\fielditem{entrysubtype}{literal}

\texttt{newspaper}（新聞）または \texttt{magazine}（商業雑誌）と書きます。

\fielditem{issue}{literal}

新聞の朝刊・夕刊です。

\fielditem{eventtitle}{literal}

学会の名前（大会名）です。\bibfield{venue} は開催地です。

\fielditem{howpublished}{literal}

雑誌では，巻・号・ページの代わりに出す文字です。コンピュータファイルでは，入手先です。

\fielditem{note}{literal}

書籍では，末尾の括弧に出ます（電子書籍のフォーマット：\texttt{Kindle版} など）。新聞では，版・地域版（\texttt{第3版，東京都版}）を書きます。

\fielditem{pages}{range}

ページです。\texttt{10--20} と書き，範囲の重複する位は自動で省かれます（\texttt{102--104} は \texttt{102-4}）。

\fielditem{url}{uri}

URL です。取得日（\bibfield{urldate}）がないと出ません。\bibfield{doi} は DOI で，取得日があるときだけ出ます（オプション \opt{doi=true} で，取得日がなくても出せます）。\bibfield{eprinttype} は，データベース名です。
\end{fieldlist}

\subsubsection{別のエントリとの関係・言語}
\begin{fieldlist}
\fielditem{crossref}{entry key}

親のエントリのキーです（\ref{sec:crossref} 節）。

\fielditem{related}{entry key}

原典のキーです。\bibfield{relatedtype} に \texttt{translationof}（訳書）または \texttt{reprintof}（再録・再刊）を書きます（\ref{sec:related} 節）。

\fielditem{langid}{identifier}

文献の言語です（\texttt{japanese}，\texttt{american} など）。和文か欧文かは，これで決まります。
\end{fieldlist}

\subsection{書き方の注意（迷いやすい点）}\label{sec:cautions}
\subsubsection{\bibfield{langid}：和文か欧文か}
和文のエントリには \texttt{langid = \{japanese\}}，欧文のエントリには \texttt{langid = \{american\}} を書きます。区切りの記号（「，」か「,」か），括弧（『』か斜体か），「編」か \texttt{ed.} か，といった書式が，これで切り替わります。

\begin{caution}
\bibfield{language} は \bibfield{langid} の代わりになりません。\bibfield{language} は文献の言語を記録するだけで，書式を切り替えません。書式を切り替えるのは \bibfield{langid} です。\bibfield{langid} がない，または文書に読み込まれていない言語のときは，欧文（\texttt{american}）として組まれます。
\end{caution}

\subsubsection{\bibfield{sortname}：和文の並び順}
文献リストは，著者名のローマ字順に並びます。漢字の読みは自動では得られないので，和文の著者・編者には，並べ替え用に \bibfield{sortname}（\texttt{Family, Given}。団体は \texttt{\{\{団体名\}\}}）を書きます。書かないと，漢字のコード順になって，欧文の文献との並びから外れ，組版時に警告（\texttt{has no sortname}）が出ます。欧文は \bibfield{sortname} が要りません。

\subsubsection{名前の書き方}
\begin{itemize}
\item 和文は \texttt{\{姓,名\}}（姓と名の間を「,」で区切る）。姓と名を区切らずに書くと，全体が姓として扱われます。
\item 欧文は \texttt{\{Family, Given\}}。名がイニシャルのときは \texttt{\{Doe, J.\}} と書きます（書いたとおりに出ます）。
\item 複数の著者は \texttt{and} でつなぎます（「,」や「，」でつながない）。他の著者を省略するときは，最後に \texttt{and others} を付けます（「ほか」／\texttt{et al.}）。
\item 団体・機関が著者のときは，二重の波括弧で囲みます（\texttt{\{\{架空総合研究所\}\}}）。囲まないと，姓と名に分けられてしまいます。
\item 姓がカタカナの外国人名は，「名・姓」の順に中黒でつながります（\texttt{\{ノヴァーク,ヤン\}} →「ヤン・ノヴァーク」）。
\end{itemize}

\subsubsection{日付とページ}
年は \bibfield{date} に書きます（\texttt{date = \{2015\}}）。BibTeX の \bibfield{year}・\bibfield{month} も読めますが，\bibfield{date} を勧めます。新聞・雑誌・ブログは年月日（\texttt{1975-03-17}）で書きます。刊行中で終了年のない文献は \texttt{1956/} です。ページの範囲は \texttt{10--20}（ハイフン2つ）で書きます。

\subsubsection{タイトルの大文字と小文字}
スタイルは，欧文タイトルの大文字小文字を変えません。書いたとおりに出ます。データ側で，ガイドの書き方（主要な語の先頭を大文字にするなど）にそろえてください。bib の中で大文字を保ちたい語は，波括弧で囲みます（\texttt{\{Tokyo\}}）。

\subsubsection{特殊文字}
\texttt{\&}，\texttt{\%}，\texttt{\#}，\texttt{\_} は，\TeX{} の特殊文字です。bib の中では \texttt{\textbackslash\&} のように，前に \texttt{\textbackslash} を付けて書きます。和文の括弧や記号は，全角のまま書けます。

\subsubsection{BibTeX の書き方との対応}
BibTeX 用に書いた bib は，biber が次の対応で読み替えるので，多くの場合そのまま使えます。

\begin{center}\small
\begin{tabular}{@{}lll@{}}
\toprule
\textbf{BibTeX} & \textbf{このスタイルで使う名前} & \textbf{備考} \\
\midrule
\texttt{journal} & \bibfield{journaltitle} & \\
\texttt{address} & \bibfield{location} & \\
\texttt{school} & \bibfield{institution} & \bibtype{thesis} \\
\texttt{year}，\texttt{month} & \bibfield{date} & \\
\texttt{howpublished} & \bibfield{howpublished} & 同じ名前 \\
\bottomrule
\end{tabular}
\end{center}

\subsubsection{\bibfield{note}：出力に出る場合}
書籍と編書論文では，\bibfield{note} が末尾に括弧で出ます（電子書籍のフォーマットを書くため）。Zotero が付ける注記（OCLC 番号など）も出力に出るので，不要なものは bib から消してください。

\subsubsection{\bibfield{url} と \bibfield{urldate}}
URL は，取得日（\bibfield{urldate}）があるときだけ出ます。書籍の書誌 URL（NDL や CiNii のページなど）を \bibfield{url} に書いておいても，取得日がなければ出力には出ません。電子媒体で利用した文献には，取得日も書きます。
